\documentclass[12pt,a4paper]{article}

\usepackage[utf8]{inputenc}
\usepackage[T1]{fontenc}
\usepackage[french]{babel}
\usepackage{geometry}
\usepackage{booktabs}
\usepackage{siunitx}
\usepackage{array}
\usepackage{xcolor}
\usepackage{titlesec}
\usepackage{fancyhdr}
\usepackage{pgfplots}

\geometry{margin=2.3cm}
\definecolor{bleufonce}{HTML}{123B63}
\definecolor{bleuclair}{HTML}{2F80C1}
\definecolor{turquoise}{HTML}{159A9C}
\definecolor{orange}{HTML}{E18A36}
\definecolor{grisclair}{HTML}{F1F5F9}
\pgfplotsset{compat=1.17}
\titleformat{\section}
  {\Large\bfseries\color{bleufonce}}
  {\thesection}{0.7em}{}
  [\vspace{0.2em}{\color{bleuclair}\titlerule}]
\titleformat{\subsection}
  {\large\bfseries\color{bleuclair}}
  {\thesubsection}{0.7em}{}
\pagestyle{fancy}
\fancyhf{}
\fancyhead[L]{\small\color{bleufonce}Benchmark des tas min}
\fancyhead[R]{\small\color{bleufonce}ERRAMI Youssef}
\fancyfoot[C]{\color{bleufonce}\thepage}
\setlength{\headheight}{15pt}
\sisetup{
  output-decimal-marker = {,},
  group-separator = {\,},
  group-minimum-digits = 4
}

\title{Comparaison de conteneurs pour l'implémentation d'un tas min}
\author{ERRAMI Youssef}
\date{30 octobre 2026}

\begin{document}
\begin{center}
  {\color{bleuclair}\rule{\textwidth}{1.2pt}}\\[0.8em]
  {\Huge\bfseries\color{bleufonce}
    Comparaison de conteneurs\\[0.2em]
    pour un tas min}\\[0.6em]
  {\Large\color{bleuclair}Comparaison des performances en C++}\\[1.2em]
  {\large\bfseries ERRAMI Youssef}\\[0.3em]
  {\color{bleufonce}30 octobre 2026}\\[0.8em]
  {\color{bleuclair}\rule{\textwidth}{1.2pt}}
\end{center}
\vspace{0.8em}

\noindent
\colorbox{grisclair}{%
  \parbox{\dimexpr\textwidth-2\fboxsep\relax}{%
    \textcolor{bleufonce}{\textbf{Comparaison.}}
    Ce rapport compare les temps d'insertion, de recherche, de suppression
    du minimum et d'affichage avec \texttt{vector}, \texttt{list} et
    \texttt{map}.%
  }%
}

\section{Environnement de test}

\begin{itemize}
  \item Ordinateur : MSI Stealth GS66 12UGS.
  \item Processeur : Intel Core i7-12700H, 14 coeurs.
  \item Carte graphique : NVIDIA GeForce RTX 3070 Ti.
  \item Mémoire vive : 16 Go DDR5.
\end{itemize}

\section{Résultats}

\subsection{Insertion}

\begin{table}[htbp]
\centering
\caption{Temps médian d'insertion (ms)}
\label{tab:insertion}
\begin{tabular}{r S[table-format=3.6] S[table-format=3.6] S[table-format=3.6]}
\toprule
{$N$} & {\texttt{vector}} & {\texttt{list}} & {\texttt{map}} \\
\midrule
100  & 0.004400 & 0.031400 & 0.037800 \\
500  & 0.015700 & 0.640700 & 0.188000 \\
1000 & 0.029100 & 2.938900 & 0.367700 \\
2000 & 0.051300 & 14.957100 & 0.875000 \\
5000 & 0.125200 & 128.559500 & 2.044300 \\
\bottomrule
\end{tabular}
\end{table}

\subsection{Recherche}

\begin{table}[htbp]
\centering
\caption{Temps médian de recherche (ms)}
\label{tab:recherche}
\begin{tabular}{r S[table-format=3.6] S[table-format=3.6] S[table-format=3.6]}
\toprule
{$N$} & {\texttt{vector}} & {\texttt{list}} & {\texttt{map}} \\
\midrule
100  & 0.001000 & 0.022500 & 0.014700 \\
500  & 0.002200 & 1.231700 & 0.132500 \\
1000 & 0.004300 & 4.685800 & 0.280400 \\
2000 & 0.008000 & 25.391200 & 0.894800 \\
5000 & 0.019300 & 361.234600 & 2.830600 \\
\bottomrule
\end{tabular}
\end{table}

\subsection{Suppression du minimum}

\begin{table}[htbp]
\centering
\caption{Temps médian de suppression répétée du minimum (ms)}
\label{tab:suppression}
\begin{tabular}{r S[table-format=3.6] S[table-format=3.6] S[table-format=3.6]}
\toprule
{$N$} & {\texttt{vector}} & {\texttt{list}} & {\texttt{map}} \\
\midrule
100  & 0.007200 & 0.039400 & 0.057400 \\
500  & 0.034000 & 1.214800 & 0.455600 \\
1000 & 0.070800 & 5.464300 & 1.124900 \\
2000 & 0.159000 & 29.574400 & 3.522300 \\
5000 & 0.451400 & 244.883200 & 10.347000 \\
\bottomrule
\end{tabular}
\end{table}

\subsection{Affichage}

\begin{table}[htbp]
\centering
\caption{Temps médian de création de la chaîne d'affichage (ms)}
\label{tab:affichage}
\begin{tabular}{r S[table-format=3.6] S[table-format=3.6] S[table-format=3.6]}
\toprule
{$N$} & {\texttt{vector}} & {\texttt{list}} & {\texttt{map}} \\
\midrule
100  & 0.026100 & 0.011900 & 0.010200 \\
500  & 0.062800 & 0.259400 & 0.057800 \\
1000 & 0.069200 & 0.966400 & 0.096600 \\
2000 & 0.147700 & 6.062300 & 0.247500 \\
5000 & 0.305800 & 50.326700 & 0.583100 \\
\bottomrule
\end{tabular}
\end{table}

\section{Visualisation graphique}

Temps en millisecondes. Les axes sont logarithmiques pour faciliter la
comparaison entre les conteneurs.

\pgfplotsset{
  benchmark axis/.style={
    width=0.92\textwidth,
    height=6.2cm,
    xmode=log,
    ymode=log,
    xmin=100,
    xmax=5000,
    xtick={100,500,1000,2000,5000},
    xticklabels={100,500,1000,2000,5000},
    xlabel={Nombre d'éléments $N$},
    ylabel={Temps médian (ms)},
    grid=both,
    major grid style={gray!30},
    minor grid style={gray!12},
    tick label style={font=\small},
    label style={font=\small\color{bleufonce}},
    legend style={
      at={(0.5,-0.25)},
      anchor=north,
      legend columns=3,
      draw=none,
      font=\small
    },
    line width=1.1pt,
    mark size=2.5pt
  }
}

\subsection{Insertion}
\begin{figure}[htbp]
\centering
\begin{tikzpicture}
\begin{axis}[benchmark axis]
\addplot[color=bleufonce,mark=*] coordinates {
  (100,0.0044) (500,0.0157) (1000,0.0291) (2000,0.0513) (5000,0.1252)
};
\addlegendentry{\texttt{vector}}
\addplot[color=turquoise,mark=square*] coordinates {
  (100,0.0314) (500,0.6407) (1000,2.9389) (2000,14.9571) (5000,128.5595)
};
\addlegendentry{\texttt{list}}
\addplot[color=orange,mark=triangle*] coordinates {
  (100,0.0378) (500,0.1880) (1000,0.3677) (2000,0.8750) (5000,2.0443)
};
\addlegendentry{\texttt{map}}
\end{axis}
\end{tikzpicture}
\caption{Temps médian d'insertion de $N$ éléments.}
\label{fig:insertion}
\end{figure}

\subsection{Recherche}
\begin{figure}[htbp]
\centering
\begin{tikzpicture}
\begin{axis}[benchmark axis]
\addplot[color=bleufonce,mark=*] coordinates {
  (100,0.0010) (500,0.0022) (1000,0.0043) (2000,0.0080) (5000,0.0193)
};
\addlegendentry{\texttt{vector}}
\addplot[color=turquoise,mark=square*] coordinates {
  (100,0.0225) (500,1.2317) (1000,4.6858) (2000,25.3912) (5000,361.2346)
};
\addlegendentry{\texttt{list}}
\addplot[color=orange,mark=triangle*] coordinates {
  (100,0.0147) (500,0.1325) (1000,0.2804) (2000,0.8948) (5000,2.8306)
};
\addlegendentry{\texttt{map}}
\end{axis}
\end{tikzpicture}
\caption{Temps médian pour jusqu'à 20 recherches de valeurs présentes.}
\label{fig:recherche}
\end{figure}

\subsection{Suppression du minimum}
\begin{figure}[htbp]
\centering
\begin{tikzpicture}
\begin{axis}[benchmark axis]
\addplot[color=bleufonce,mark=*] coordinates {
  (100,0.0072) (500,0.0340) (1000,0.0708) (2000,0.1590) (5000,0.4514)
};
\addlegendentry{\texttt{vector}}
\addplot[color=turquoise,mark=square*] coordinates {
  (100,0.0394) (500,1.2148) (1000,5.4643) (2000,29.5744) (5000,244.8832)
};
\addlegendentry{\texttt{list}}
\addplot[color=orange,mark=triangle*] coordinates {
  (100,0.0574) (500,0.4556) (1000,1.1249) (2000,3.5223) (5000,10.3470)
};
\addlegendentry{\texttt{map}}
\end{axis}
\end{tikzpicture}
\caption{Temps médian de suppression répétée du minimum.}
\label{fig:suppression}
\end{figure}

\subsection{Affichage}
\begin{figure}[htbp]
\centering
\begin{tikzpicture}
\begin{axis}[benchmark axis]
\addplot[color=bleufonce,mark=*] coordinates {
  (100,0.0261) (500,0.0628) (1000,0.0692) (2000,0.1477) (5000,0.3058)
};
\addlegendentry{\texttt{vector}}
\addplot[color=turquoise,mark=square*] coordinates {
  (100,0.0119) (500,0.2594) (1000,0.9664) (2000,6.0623) (5000,50.3267)
};
\addlegendentry{\texttt{list}}
\addplot[color=orange,mark=triangle*] coordinates {
  (100,0.0102) (500,0.0578) (1000,0.0966) (2000,0.2475) (5000,0.5831)
};
\addlegendentry{\texttt{map}}
\end{axis}
\end{tikzpicture}
\caption{Temps médian de préparation de l'affichage du tas.}
\label{fig:affichage}
\end{figure}

\section{Comparaison}

\texttt{vector} est le plus rapide pour l'insertion, la recherche et la
suppression du minimum. \texttt{map} arrive généralement en deuxième position,
tandis que \texttt{list} est la plus lente pour les opérations du tas. Ces
résultats s'expliquent par les accès par indice nécessaires au tas binaire.

\end{document}
